\paragraph{Walkthrough.}
The initial graph uses a bounded propose--critique loop. It performs
\texttt{num\_rounds} reviews and one final planner pass, then a caller places
\texttt{current\_plan} into the separately compiled adaptive graph. There,
\texttt{execute\_step} reuses \texttt{deep\_research.run\_step}: it dispatches the
current step to the engineer's bounded self-debug graph or the researcher graph,
records a summary and structured outcome, checkpoints the run, and increments
\texttt{step\_index}.

The adaptive reviewer is called after every executed step. Given the main task,
latest summary and outcome, and current unexecuted tail, its structured
\texttt{needs\_adaptation} field triggers \texttt{replan\_tail}. The replanner may
replace every not-yet-executed step, including their number, contents and assigned
agents. \texttt{merge\_plan} holds the completed prefix
\texttt{plan[:step\_index-1]} and its outputs fixed, then appends the replacement
tail. Otherwise execution advances directly. Routing stops after a failed step or
when \texttt{step\_index > len(plan)}. This guarantees termination only while the
remaining plan is finite and is not repeatedly extended. Unlike the initial
planning and inner self-debug loops, the adaptive loop has no total-step or
revision bound, so the implementation does not provide an unconditional
termination guarantee.
